\documentclass[11pt]{article}
\usepackage[margin=1in]{geometry}
\usepackage{amsmath, amssymb, amsthm}
\usepackage{algorithm}
\usepackage{algorithmic}
\usepackage{color}

\title{Analysis of Randomized Coordinate Descent under Polyak-\L{}ojasiewicz Condition}
\author{Optimization Theory Analysis}
\date{}

\begin{document}
\maketitle

\section*{Algorithm Name}
Randomized Coordinate Descent (RCD)

\section*{Mathematical Formulation}
The algorithm aims to minimize a continuously differentiable objective function $f: \mathbb{R}^d \to \mathbb{R}$. At each iteration $t$, the algorithm selects a coordinate index $i_t \in \{1, 2, \dots, d\}$ uniformly at random. It then computes the partial derivative of the objective function with respect to the chosen coordinate, $\nabla_{i_t} f(w_t)$, and updates the full parameter vector $w_t$ by taking a step in the direction of the corresponding standard basis vector $e_{i_t}$.

The formal update rule is:
\begin{equation}
    w_{t+1} = w_t - \eta \nabla_{i_t} f(w_t) e_{i_t}
\end{equation}
where:
\begin{itemize}
    \item $w_t \in \mathbb{R}^d$ is the parameter vector at iteration $t$.
    \item $i_t \sim \text{Uniform}(1, d)$ is the randomly chosen coordinate index.
    \item $\nabla_{i_t} f(w_t) = \frac{\partial f}{\partial w^{(i_t)}}(w_t)$ is the coordinate-wise gradient (a scalar).
    \item $e_{i_t} \in \mathbb{R}^d$ is the standard basis vector with $1$ at position $i_t$ and $0$ elsewhere.
    \item $\eta > 0$ is the learning rate (step size).
\end{itemize}

\section*{Pseudocode}
\begin{algorithm}[H]
\caption{Randomized Coordinate Descent (RCD)}
\begin{algorithmic}[1]
\REQUIRE Objective function $f$, step size $\eta > 0$, total iterations $T$, initial parameters $w_0 \in \mathbb{R}^d$
\FOR{$t = 0$ to $T-1$}
    \STATE Sample coordinate index $i_t$ uniformly at random from $\{1, 2, \dots, d\}$
    \STATE Compute the partial derivative: $g_t = \nabla_{i_t} f(w_t)$
    \STATE Form the sparse update vector: $\Delta_t = g_t e_{i_t}$
    \STATE Update the parameter vector: $w_{t+1} = w_t - \eta \Delta_t$
\ENDFOR
\RETURN $w_T$
\end{algorithmic}
\end{algorithm}

\section*{Dry Run Example}
We apply the algorithm to the 1-dimensional objective function $f(w) = (w-3)^2$ with $w \in \mathbb{R}$. 
Since $d=1$, the randomly selected coordinate is always $i_t = 1$.
The gradient is $\nabla f(w) = 2(w-3)$. We initialize $w_0 = 0$ and use a step size $\eta = 0.1$.

\textbf{Initialization:}
$w_0 = 0$, $f(w_0) = (-3)^2 = 9$.

\textbf{Iteration 1 ($t=0$):}
\begin{itemize}
    \item Gradient: $g_0 = \nabla f(w_0) = 2(0 - 3) = -6$
    \item Update: $w_1 = w_0 - \eta g_0 = 0 - 0.1(-6) = 0.6$
    \item Objective: $f(w_1) = (0.6 - 3)^2 = (-2.4)^2 = 5.76$
\end{itemize}

\textbf{Iteration 2 ($t=1$):}
\begin{itemize}
    \item Gradient: $g_1 = \nabla f(w_1) = 2(0.6 - 3) = -4.8$
    \item Update: $w_2 = w_1 - \eta g_1 = 0.6 - 0.1(-4.8) = 1.08$
    \item Objective: $f(w_2) = (1.08 - 3)^2 = (-1.92)^2 = 3.6864$
\end{itemize}

\textbf{Iteration 3 ($t=2$):}
\begin{itemize}
    \item Gradient: $g_2 = \nabla f(w_2) = 2(1.08 - 3) = -3.84$
    \item Update: $w_3 = w_2 - \eta g_2 = 1.08 - 0.1(-3.84) = 1.464$
    \item Objective: $f(w_3) = (1.464 - 3)^2 = (-1.536)^2 = 2.359296$
\end{itemize}

\section*{Assumptions}
To analyze the convergence of this algorithm, we make two primary assumptions:

\textbf{Assumption 1 (Coordinate-wise Smoothness):} 
The function $f$ has coordinate-wise Lipschitz continuous gradients. There exists a constant $L > 0$ such that for any $w \in \mathbb{R}^d$, $h \in \mathbb{R}$, and any coordinate $i \in \{1, \dots, d\}$:
\begin{equation}
    f(w + h e_i) \leq f(w) + h \nabla_i f(w) + \frac{L}{2} h^2
\end{equation}

\textbf{Assumption 2 (Polyak-\L{}ojasiewicz (PL) Condition):}
The function $f$ satisfies the PL inequality with constant $\mu > 0$. For all $w \in \mathbb{R}^d$:
\begin{equation}
    \frac{1}{2} \|\nabla f(w)\|^2 \geq \mu (f(w) - f^*)
\end{equation}
where $f^* = \min_w f(w)$ is the global minimum of the function. Note that the PL condition does not require the function to be convex, yet it guarantees that every stationary point is a global minimum.

\section*{Main Convergence Theorem}
\textbf{Theorem 1.} Under Assumptions 1 and 2, if we run Randomized Coordinate Descent with a constant step size $\eta = \frac{1}{L}$, the expected suboptimality at iteration $t$ is bounded by:
\begin{equation}
    \mathbb{E}[f(w_t) - f^*] \leq \left(1 - \frac{\mu}{dL}\right)^t (f(w_0) - f^*)
\end{equation}
This establishes a global linear convergence rate to the optimal value $f^*$.

\section*{Theoretical Properties}
\begin{itemize}
    \item \textbf{Convergence Rate:} The convergence rate is $\mathcal{O}(\rho^t)$ where $\rho = 1 - \frac{\mu}{dL}$. This represents linear convergence, which is highly efficient.
    \item \textbf{Comparison to Gradient Descent (GD):} Standard GD under $L_{full}$-smoothness and PL condition achieves a rate of $(1 - \mu/L_{full})^t$. Since $L \leq L_{full} \leq dL$, RCD often requires $d$ times more iterations than GD. However, each RCD iteration is $d$ times cheaper to compute, making the overall computational complexity comparable or strictly better if $L \ll L_{full}$.
    \item \textbf{Hyperparameter Sensitivity:} The algorithm is sensitive to the step size $\eta$. If $\eta > \frac{2}{L}$, the algorithm may diverge. The optimal theoretical constant step size is $\eta = \frac{1}{L}$.
\end{itemize}

\section*{Discussion}
The PL condition is a relaxed assumption compared to strong convexity. It allows the function to be non-convex (e.g., in certain neural network formulations or matrix factorization problems) while still enabling global linear convergence. RCD is particularly advantageous when the cost of computing a single partial derivative is significantly less than computing the full gradient.

\section*{Preliminaries (Definitions and Lemmas)}
\textbf{Lemma 1 (Descent Lemma for Coordinate Update):}
Given Assumption 1 and the update rule $w_{t+1} = w_t - \eta \nabla_{i_t} f(w_t) e_{i_t}$, we have:
\begin{equation}
    f(w_{t+1}) \leq f(w_t) - \eta (\nabla_{i_t} f(w_t))^2 + \frac{L}{2} \eta^2 (\nabla_{i_t} f(w_t))^2
\end{equation}
\textit{Proof:} Directly follows by substituting $h = -\eta \nabla_{i_t} f(w_t)$ and $i = i_t$ into Assumption 1.

\section*{Main Proof (Step-by-Step Derivation)}

We now provide the complete, step-by-step derivation of the convergence rate. Let $\mathbb{E}_{i_t}$ denote the conditional expectation with respect to the random choice of the coordinate $i_t$ at step $t$, given all previous history (i.e., given $w_t$). Let $\mathbb{E}$ denote the full expectation over all randomness.

\begin{align}
    \text{[Step 1]} \quad & f(w_{t+1}) \leq f(w_t) - \eta \nabla_{i_t} f(w_t)^2 + \frac{L}{2} \eta^2 \nabla_{i_t} f(w_t)^2 \\
    & \text{(Applying the coordinate-wise smoothness assumption with $h = -\eta \nabla_{i_t} f(w_t)$)} \nonumber \\[10pt]
    \text{[Step 2]} \quad & f(w_{t+1}) \leq f(w_t) - \left( \eta - \frac{L}{2} \eta^2 \right) (\nabla_{i_t} f(w_t))^2 \\
    & \text{(Factoring out the squared partial derivative term)} \nonumber \\[10pt]
    \text{[Step 3]} \quad & f(w_{t+1}) \leq f(w_t) - \frac{1}{2L} (\nabla_{i_t} f(w_t))^2 \\
    & \text{(Substituting the optimal step size $\eta = \frac{1}{L}$, so $\eta - \frac{L}{2}\eta^2 = \frac{1}{L} - \frac{1}{2L} = \frac{1}{2L}$)} \nonumber \\[10pt]
    \text{[Step 4]} \quad & \mathbb{E}_{i_t}[f(w_{t+1})] \leq \mathbb{E}_{i_t} \left[ f(w_t) - \frac{1}{2L} (\nabla_{i_t} f(w_t))^2 \right] \\
    & \text{(Taking the conditional expectation given $w_t$ on both sides)} \nonumber \\[10pt]
    \text{[Step 5]} \quad & \mathbb{E}_{i_t} [(\nabla_{i_t} f(w_t))^2] = \sum_{i=1}^d \frac{1}{d} (\nabla_i f(w_t))^2 = \frac{1}{d} \|\nabla f(w_t)\|^2 \\
    & \text{(Evaluating the expectation since $i_t$ is chosen uniformly from $\{1, \dots, d\}$)} \nonumber \\[10pt]
    \text{[Step 6]} \quad & \mathbb{E}_{i_t}[f(w_{t+1})] \leq f(w_t) - \frac{1}{2dL} \|\nabla f(w_t)\|^2 \\
    & \text{(Substituting the result of [Step 5] into [Step 4])} \nonumber \\[10pt]
    \text{[Step 7]} \quad & \mathbb{E}_{i_t}[f(w_{t+1})] - f^* \leq f(w_t) - f^* - \frac{1}{2dL} \|\nabla f(w_t)\|^2 \\
    & \text{(Subtracting the global minimum $f^*$ from both sides)} \nonumber \\[10pt]
    \text{[Step 8]} \quad & - \|\nabla f(w_t)\|^2 \leq - 2\mu (f(w_t) - f^*) \\
    & \text{(Rearranging the Polyak-\L{}ojasiewicz (PL) condition: $\frac{1}{2}\|\nabla f(w)\|^2 \geq \mu(f(w)-f^*)$)} \nonumber \\[10pt]
    \text{[Step 9]} \quad & \mathbb{E}_{i_t}[f(w_{t+1})] - f^* \leq f(w_t) - f^* - \frac{1}{2dL} \left( 2\mu (f(w_t) - f^*) \right) \\
    & \text{(Applying the inequality from [Step 8] into the bound from [Step 7])} \nonumber \\[10pt]
    \text{[Step 10]} \quad & \mathbb{E}_{i_t}[f(w_{t+1})] - f^* \leq \left( 1 - \frac{\mu}{dL} \right) (f(w_t) - f^*) \\
    & \text{(Factoring out the suboptimality term $(f(w_t) - f^*)$)} \nonumber \\[10pt]
    \text{[Step 11]} \quad & \mathbb{E}[f(w_{t+1}) - f^*] \leq \left( 1 - \frac{\mu}{dL} \right) \mathbb{E}[f(w_t) - f^*] \\
    & \text{(Taking the full expectation $\mathbb{E}$ over all random variables up to step $t$)} \nonumber \\[10pt]
    \text{[Step 12]} \quad & \mathbb{E}[f(w_t) - f^*] \leq \left( 1 - \frac{\mu}{dL} \right)^t (f(w_0) - f^*) \\
    & \text{(Unrolling the recurrence relation from $t$ down to $0$)} \nonumber
\end{align}

\section*{Rate Derivation (With Constants)}
From [Step 12], we have established the linear convergence rate. Let us define the condition number with respect to the coordinate-wise smoothness as $\kappa_{coord} = \frac{L}{\mu}$. 
The contraction factor is $\rho = 1 - \frac{1}{d \kappa_{coord}}$.

To find the number of iterations $T$ required to achieve an expected $\epsilon$-suboptimality, i.e., $\mathbb{E}[f(w_T) - f^*] \leq \epsilon$, we set:
\begin{equation}
    \left( 1 - \frac{\mu}{dL} \right)^T (f(w_0) - f^*) \leq \epsilon
\end{equation}
Taking the natural logarithm of both sides:
\begin{equation}
    T \ln\left( 1 - \frac{\mu}{dL} \right) \leq \ln\left( \frac{\epsilon}{f(w_0) - f^*} \right)
\end{equation}
Using the inequality $\ln(1-x) \leq -x$ for $x \in (0, 1)$:
\begin{equation}
    -T \frac{\mu}{dL} \leq \ln\left( \frac{\epsilon}{f(w_0) - f^*} \right) \implies T \geq \frac{dL}{\mu} \ln\left( \frac{f(w_0) - f^*}{\epsilon} \right)
\end{equation}
Thus, the iteration complexity is $\mathcal{O}\left( d \kappa_{coord} \log\frac{1}{\epsilon} \right)$.

\section*{Special Cases}
\textbf{1. Strongly Convex Functions:}
If the function $f$ is $\mu$-strongly convex, it automatically satisfies the PL condition with the same constant $\mu$. Therefore, the derived linear convergence rate $\mathcal{O}\left((1 - \frac{\mu}{dL})^t\right)$ strictly holds for any $\mu$-strongly convex function that is coordinate-wise $L$-smooth.

\textbf{2. Convex Quadratic Functions:}
Consider $f(w) = \frac{1}{2} w^\top Q w - c^\top w$ where $Q \in \mathbb{R}^{d \times d}$ is a positive definite matrix with eigenvalues $0 < \lambda_{\min} \leq \dots \leq \lambda_{\max}$. 
Here, the coordinate-wise smoothness constant $L$ is bounded by the maximum diagonal entry of $Q$, i.e., $L = \max_i Q_{ii}$. The PL constant is $\mu = \lambda_{\min}$.
The convergence rate becomes $\left( 1 - \frac{\lambda_{\min}}{d \max_i Q_{ii}} \right)^t$. 
Because $\max_i Q_{ii} \leq \lambda_{\max}$, RCD on quadratics can be significantly faster than full gradient descent when the diagonal elements of $Q$ are much smaller than its maximum eigenvalue.

\end{document}